\subsection{Quantum Annealing}\label{sec:quantum-annealing}

In the previous sections, we established that a QUBO problem can be transformed into an Ising Hamiltonian whose ground state corresponds to the optimal solution of the original optimization problem.

\highspace
However, knowing the Hamiltonian is not enough. We still need a mechanism capable of driving the quantum system toward its ground state.

\highspace
\textbf{Quantum Annealing} is a \textbf{computational paradigm designed} precisely \textbf{for this purpose}. Instead of searching directly among all possible solutions, the system is initialized in the ground state of a simple Hamiltonian and then slowly evolved toward the Hamiltonian that encodes the optimization problem.

\highspace
If this evolution is performed sufficiently slowly, the system tends to remain in its instantaneous ground state throughout the process. As a consequence, the final state of the system approximates the ground state of the problem Hamiltonian, which corresponds to the optimal solution of the original QUBO problem.

\highspace
If these concepts seem too abstract, don't worry! We will cover the specifics of quantum annealing in the following sections.